\begin{problem}[Finite reduced algebras are finite sets of points]
Let $k$ be algebraically closed and let $A$ be a finite-dimensional commutative
$k$-algebra.  Prove that if $A$ is reduced, then
\[
A\cong k^r
\]
for some $r$.  Interpret the result geometrically.
\end{problem}
\paragraph{Why this problem matters.}
It isolates the exact role of nilpotents in zero-dimensional geometry: without them, a
finite affine scheme over an algebraically closed field is just a finite set of points.
\paragraph{Useful ideas before starting.}
Finite-dimensional algebras are Artinian.  Use their finitely many maximal ideals and the
Chinese remainder theorem.
\paragraph{Guided hints.}
In a reduced Artinian ring, the nilradical is zero and equals the intersection of all
maximal ideals.
\begin{solution}
Let $\mathfrak m_1,\dots,\mathfrak m_r$ be the maximal ideals.  In an Artinian ring every
prime is maximal and the nilradical is their intersection.  Reducedness gives
\[
\bigcap_i\mathfrak m_i=0.
\]
Distinct maximal ideals are comaximal, so CRT yields
\[
A\cong\prod_i A/\mathfrak m_i.
\]
Each factor is a finite field extension of $k$, hence equals $k$ because $k$ is algebraically
closed.  Thus $A\cong k^r$.  Its spectrum is a discrete set of $r$ reduced points.
\end{solution}
\paragraph{What happened mathematically.}
Nilpotents are the only source of nontrivial infinitesimal structure in finite affine
geometry over an algebraically closed field.
\paragraph{Extensions.}
Drop algebraic closedness and identify the factors as finite field extensions of $k$.
\paragraph{Provenance.}
New synthesis problem using CRT material already developed in the series.
